\section{Related Work}
\label{sec:related}

\subsection{Domain Generalization for Hyperspectral Image Classification}
Domain generalization (DG) for hyperspectral image (HSI) classification aims to train a model on one or more source domains such that it generalizes directly to unseen target domains without target data access during training~\cite{wang2022domain}. In remote sensing, geographical domain shifts—caused by variations in atmospheric conditions, sensor specifications, and lighting angles—severely degrade the performance of traditional classification models. To address this issue, research trends have focused on expanding the model's generalization boundaries through data augmentation, domain-invariant feature extraction, or progressive generation. Data augmentation and style randomization frameworks, such as SDENet~\cite{sdenet2022} and LLURNet~\cite{llurnet2023}, focus on simulating diverse target-domain distributions by perturbing visual and spectral statistics. While these methods successfully broaden the visual search space, random perturbations applied directly to feature maps may not explicitly preserve the semantic layout or class structural information. Alternatively, domain-invariant representation learning and adversarial methods, such as FDGNet~\cite{fdgnet2022} and S2AMSNet~\cite{s2amsnet2022}, utilize discrepancy metrics or adversarial classifiers to isolate domain-specific variations from target-invariant features. While effective at aligning global statistics, these adversarial constraints may introduce semantic inconsistencies in the localized spatial-spectral representations. Finally, progressive generation methods, such as PMGDG~\cite{li2022progressive}, decompose the target domain simulation into sequential band generation steps. Although progressive schemes make high-dimensional generation tractable, the style perturbations applied at intermediate stages can limit cross-domain semantic consistency. Consequently, existing DG methods primarily rely on visual feature learning and often lack explicit semantic guidance to stabilize representation mapping during cross-domain generation.

\subsection{Vision-Language Learning for Hyperspectral Image Classification}
Vision-language learning has emerged as a powerful paradigm for cross-modal representation learning, spearheaded by foundation models like CLIP~\cite{radford2021learning} and ALIGN. By aligning visual patches and textual descriptions in a shared latent space via contrastive objectives, these models capture rich, domain-agnostic semantic concepts. In remote sensing, researchers have adapted vision-language concepts to scene classification and target detection to leverage the zero-shot generalization capabilities of linguistic priors. Specifically for HSI domain generalization, models like EHSNet~\cite{ehsnet2023} and asymmetric contrastive bridging models like ACB~\cite{acb2023} utilize class-level text descriptions to align visual feature maps in a text-guided latent space, improving the robust classification of out-of-distribution ground objects. While these approaches demonstrate the value of cross-modal alignment, adapting natural-image vision-language models to HSI remains highly challenging. First, standard vision-language models are designed for three-channel RGB images and struggle to extract the continuous spatial-spectral signatures of multi-band HSI cubes. Second, designing effective, context-rich textual prompts remains complex, often relying on manual annotation. Most importantly, existing visual-language models do not dynamically utilize language features to guide spatial-spectral generation. These limitations leave a critical gap: existing methods struggle to adapt vision-language representations dynamically to guide spatial-spectral synthesis, leading to insufficient semantic conditioning during feature generation.

\subsection{Progressive Generation and Semantic Alignment}
Progressive generation strategies decompose the modeling of complex, high-dimensional target distributions into a cascade of lower-dimensional generation steps. For HSI domain generalization, progressive domain expansion and step-wise spectral downsampling cascades make the optimization of high-dimensional HSI patches tractable, simulating realistic target-domain variations. Existing progressive generation models primarily focus on expanding visual diversity and learning domain-invariant representations through feature generation. By partitioning the high-dimensional HSI bands into progressive spectral steps, these methods reduce the complexity of the feature translation task. Conversely, recent vision-language methods mainly focus on improving semantic representation learning through cross-modal image-text alignment in a shared latent space.

However, these two research directions are largely studied independently. Integrating semantic guidance throughout the progressive generation process to improve domain-invariant feature generation remains relatively underexplored in HSI domain generalization. Incorporating semantic guidance has the potential to preserve class semantics during the progressive generation steps, improving the semantic consistency and transferability of generated feature representations. To address these limitations, we present a unified framework in the next section that combines progressive generation, semantic guidance, vision-language alignment, and adaptive semantic conditioning.
